\documentclass[11pt,a4paper]{article}
\usepackage[utf8]{inputenc}
\usepackage[T1]{fontenc}
\usepackage[english]{babel}
\usepackage{lmodern}
\usepackage{geometry}
\usepackage{xcolor}
\usepackage{tocloft}
\usepackage{hyperref}
\usepackage{enumitem}
\usepackage{parskip}
\usepackage{booktabs}
\usepackage{array}
\usepackage{longtable}
\usepackage{fancyhdr}
\usepackage{lastpage}
\usepackage{listings}
\usepackage{titlesec}
\geometry{margin=2.15cm}
\setlength{\headheight}{14pt}
\setlength{\emergencystretch}{3em}
\setlist{itemsep=3pt,parsep=1pt,topsep=4pt}
\definecolor{titleblue}{HTML}{123B5D}
\definecolor{softblue}{HTML}{EDF5FB}
\definecolor{softgray}{HTML}{F5F7FA}
\definecolor{accent}{HTML}{0F6E8C}
\hypersetup{
  colorlinks=true,linkcolor=titleblue,urlcolor=accent,
  pdftitle={Scientific Data Analysis v1.8: English publication edition},
  pdfauthor={interemi},
  pdfsubject={Backend contracts for app integration}
}
\titleformat{\section}{\Large\bfseries\color{titleblue}}{\thesection.}{0.6em}{}
\titleformat{\subsection}{\large\bfseries\color{accent}}{\thesubsection.}{0.6em}{}
\setlength{\cftsecnumwidth}{2.8em}
\tocloftpagestyle{fancy}
\lstdefinestyle{promptstyle}{
  basicstyle=\ttfamily\small,backgroundcolor=\color{softgray},
  frame=single,rulecolor=\color{softblue},breaklines=true,columns=fullflexible
}
\pagestyle{fancy}
\fancyhf{}
\fancyhead[L]{Scientific Data Analysis}
\fancyhead[R]{v1.8 guide: English edition}
\fancyfoot[C]{\thepage\ of \pageref{LastPage}}
\newcommand{\code}[1]{\texttt{\detokenize{#1}}}

\begin{document}
\pagenumbering{gobble}
\begin{titlepage}
\centering
\vspace*{1.0cm}
{\Huge\bfseries\color{titleblue} Integration Guide\\[0.15cm]
Scientific Data Analysis v1.8\par}
\vspace{0.65cm}
{\Large Backend contracts for app integration\par}
\vspace{0.8cm}
{\large App invocation, run bundles, dry-run routing, and explicit limits\par}
\vspace{0.8cm}
{\large English publication edition: September 23, 2026\par}
\vspace{0.5cm}
\begin{minipage}{0.93\textwidth}
\small\raggedright
Translation of the preserved private v1.8 guide at commit \code{994eaea}.
Original TeX SHA-256:
\nolinkurl{41c5839931276961b04a6377cd247591fcf30940f120361b9b95a9b8583efdcf}.
Private source:
\nolinkurl{skills/scientific-data-maintainer/.cache/active_references/guides/guia_scientific_data_analysis_v1_8.tex}.

This guide records historical backend preparation, not current verification
or installed-family promotion. Its v1.8 scope does not modify or synchronize
Scientific Workbench. Historical commands are illustrative: use fresh output
directories when reproducing them without overwriting earlier evidence.
The example task prompt is translated into English; script names, options,
and contract identifiers are preserved. The original remains unchanged privately.
\end{minipage}
\vfill
{\large Prepared with Codex\par}
{\large Copyright 2026 interemi\par}
\end{titlepage}

\clearpage
\pagenumbering{roman}
\begin{abstract}
\code{scientific-data-analysis} v1.8 is a contract release. It prepares the
editable skill as a more predictable backend for Scientific Workbench,
without claiming that the app is already synchronized to v2.0. The backend
comes first: parseable JSON, artifact types, run bundles, dry-run routing,
an app-readiness matrix, app-like regressions, and explicit separation of
normal, expert, optional, legacy, maintainer, and unexposed modes.

v1.7 asked whether the skill was useful in real settings. v1.8 asks whether
an app can invoke it predictably, read the result, present artifacts, and
explain the next step while preserving original inputs.
\end{abstract}
\tableofcontents
\newpage
\pagenumbering{arabic}

\section{Executive overview}
v1.8 prepares the skill as a Scientific Workbench backend. It does not develop
or synchronize the app, move it to v2.0, or promise that the app already
consumes every new field. The backend changes are:
\begin{itemize}[leftmargin=1.5em]
\item A v1.8 JSON contract for app-ready outputs.
\item An initial artifact type taxonomy.
\item A stable run bundle format.
\item An app-readiness matrix for all 59 public capabilities.
\item A deterministic dry-run router.
\item App-like regressions for general, scientific, astronomy, optional, and legacy routes.
\item Integration guidance that reserves major app synchronization for v2.0.
\end{itemize}
\textbf{Product intent:} each capability should state whether it applies,
warns, blocks, or fails cleanly, which artifacts it produced, and what a
person should do next. v1.8 does not make every capability general-purpose.

\section{From v1.7 to v1.8}
\begin{longtable}{p{0.25\linewidth}p{0.64\linewidth}}
\toprule
\textbf{Area} & \textbf{v1.8 change} \\
\midrule
Contract & A JSON envelope with \code{contract_version}, \code{app_status}, \code{typed_artifacts}, \code{original_modified}, \code{warnings}, \code{errors}, \code{next_actions}, and \code{app_hints}. \\
App-readiness & All 59 capabilities are classified as \code{app_ready}, \code{app_ready_partial}, \code{cli_only}, \code{blocked_optional}, \code{maintainer_only}, or \code{not_applicable_to_app}. \\
Run bundles & A stable run folder with \code{summary.json}, \code{manifest.json}, logs, artifacts, previews, reports, tables, and \code{next_steps.md}. \\
Router & \code{scientific_workflow_router.py} is a conservative dry-run planner. It recommends or rejects routes without executing destructive actions. \\
App-like regressions & General, scientific, astronomy, optional, and legacy families and bundles are exercised as a local app would invoke them. \\
Integration & Guidance explains how Scientific Workbench can use the contracts later, without claiming v2.0 integration. \\
\bottomrule
\end{longtable}

\section{What remains unchanged}
\begin{itemize}[leftmargin=1.5em]
\item The installed skill remains at its applicable stable version until final synchronization.
\item Scientific Workbench is not edited in this phase.
\item Domain-specific capabilities retain their domain limits.
\item Optional routes do not become core dependencies.
\item Legacy routes are not ordinary user actions.
\item Original user inputs remain read-only by default.
\item v1.8 builds on the real-world foundation of v1.7.
\end{itemize}

\section{The v1.8 JSON contract}
Primary reference: \path{references/v1-8-app-ready-contract.md}.

Minimum fields expected in app-ready output:
\begin{longtable}{>{\raggedright\arraybackslash}p{0.36\linewidth}p{0.53\linewidth}}
\toprule
\textbf{Field} & \textbf{Use} \\
\midrule
\code{contract_version} & \code{1.8} for an app-ready envelope. \\
\code{tool} & Stable script or route name, retained for compatibility. \\
\code{capability_id} / \code{capability_label} & Registry identity where applicable. \\
\code{status} / \code{app_status} & Canonical or normalized status. \\
\code{command} & \code{argv}, \code{cwd}, and redaction. Must never expose secrets. \\
\code{inputs} & Inspected paths or copies used. \\
\code{outputs} & Expected run products. \\
\code{artifacts} / \code{typed_artifacts} & Typed artifacts for the UI. \\
\code{warnings} & Useful risks that do not invalidate the entire result. \\
\code{errors} & Clean failures or blocks, without raw tracebacks. \\
\code{qa} & Quality findings and metrics. \\
\code{provenance} & Original-input policy, environment, and source. \\
\code{next_actions} & A person's next step. \\
\code{original_modified} & \code{false} for normal routes. \\
\code{app_hints} & Short summary, visual severity, tags, and suggested previews. \\
\bottomrule
\end{longtable}

\subsection{Canonical statuses}
\begin{longtable}{p{0.29\linewidth}p{0.60\linewidth}}
\toprule
\textbf{Status} & \textbf{Meaning} \\
\midrule
\code{PASS} & The run finished and its result is usable. \\
\code{WARNING} & Useful output requires human review. \\
\code{BLOCKED_CONTROLADO} & A backend, input, platform, or safety condition prevents execution; the block is clean and actionable. \\
\code{FAIL} & Clean failure without products that appear valid. \\
\code{ROTO} & Audit classification for raw tracebacks, false success, misleading artifacts, or changed originals. It should not be a public tool's normal status. \\
\bottomrule
\end{longtable}

\section{Artifact types}
Initial taxonomy:
\begin{longtable}{p{0.29\linewidth}p{0.60\linewidth}}
\toprule
\textbf{Artifact type} & \textbf{Meaning} \\
\midrule
\code{summary_json} & App-parseable summary. \\
\code{manifest_json} & Provenance: inputs, outputs, command, and environment. \\
\code{report_md} & Human-readable Markdown report. \\
\code{preview_png} & Static visual preview. \\
\code{table_csv} & Derived table. \\
\code{notebook_ipynb} & Generated, inspected, or copy-executed notebook. \\
\code{log_txt} & stdout, stderr, command log, or transcript. \\
\code{qa_report} & Dedicated QA report. \\
\code{handoff_bundle} & Folder or package for continuation or delivery. \\
\code{unknown} & Existing artifact without a more specific type. \\
\bottomrule
\end{longtable}

\section{Run bundles}
Reference: \path{references/v1-8-run-bundle-contract.md}.
Target layout:
\begin{lstlisting}[style=promptstyle]
run/
  manifest.json
  summary.json
  stdout.txt
  stderr.txt
  command.txt
  artifacts/
  previews/
  reports/
  tables/
  logs/
  next_steps.md
\end{lstlisting}
This contract lets Scientific Workbench discover results without fragile
heuristics; it does not require rewriting every script in v1.8. Priority
general routes already have focused bundle or wrapper coverage. Scientific
and legacy routes can migrate gradually.

\section{Dry-run router}
References: \path{references/v1-8-workflow-router.md} and
\path{scripts/scientific_workflow_router.py}.
\begin{longtable}{p{0.20\linewidth}p{0.69\linewidth}}
\toprule
\textbf{Mode} & \textbf{Use} \\
\midrule
\code{inspect} & Inventories the input and proposes candidate routes. \\
\code{plan} & Returns dry-run steps and suggested commands without executing them. \\
\code{explain} & Explains why a capability applies, does not apply, or is unknown. \\
\bottomrule
\end{longtable}
Example:
\begin{lstlisting}[style=promptstyle]
python scripts/scientific_workflow_router.py plan input.csv \
  --task "profile an anonymized professional table" \
  --summary-json run/router_summary.json \
  --manifest-json run/router_manifest.json
\end{lstlisting}
The router does not run downstream tools or modify inputs. It should reject
astronomy or legacy routes for general input.

\newpage
\section{Capability app-readiness}
Reference: \path{references/v1-8-app-readiness-matrix.md}.
Recorded v1.8 counts:
\begin{longtable}{p{0.34\linewidth}r p{0.44\linewidth}}
\toprule
\textbf{Class} & \textbf{Total} & \textbf{Use} \\
\midrule
\code{app_ready} & 17 & App invocation with sufficiently clear output. \\
\code{app_ready_partial} & 18 & Useful, but needs a wrapper, typing, documentation, or UX work. \\
\code{blocked_optional} & 7 & Optional backend or platform: preflight panel, not an ordinary action. \\
\code{cli_only} & 9 & Valid CLI route, but narrow or expert-oriented for a normal app flow. \\
\code{maintainer_only} & 6 & Release, registry, smoke, or maintenance only. \\
\code{not_applicable_to_app} & 2 & Unexposed unless a clear expert context applies. \\
\bottomrule
\end{longtable}

\subsection{Recommended exposure}
\begin{itemize}[leftmargin=1.5em]
\item \textbf{Normal}: tables, documents, containers, mixed folders, time series, physical QA, and noise budgets.
\item \textbf{Expert}: FITS, WCS, astrometry, sky crossmatching, RGB/FITS, RV, photometric calibration, spectroscopy, and advanced presentations.
\item \textbf{Optional panel}: STILTS/TOPCAT, APT, DuckDB, iWork, Keynote, and TEAREDUCE.
\item \textbf{Legacy expert}: IRAF, fxcor, iSTARMOD, and narrow coursework routes.
\item \textbf{Maintainer}: smoke, synchronization, registry, capability probes, and gates.
\end{itemize}

\section{Testing a capability as an app would}
Minimum pattern:
\begin{enumerate}[leftmargin=1.5em]
\item Copy or synthesize input in a fresh \code{tmp/} run.
\item Use the dry-run router when the case is ambiguous.
\item Run the capability with \code{--summary-json} and, where available, \code{--manifest-json} or \code{--run-dir}.
\item Verify that \code{summary.json} parses.
\item Check \code{contract_version}, \code{app_status}, \code{typed_artifacts}, \code{original_modified=false}, \code{warnings}, \code{errors}, and \code{next_actions}.
\item Confirm that the original input has not changed.
\item Classify output as \code{PASS}, \code{WARNING}, \code{BLOCKED_CONTROLADO}, \code{FAIL}, or audit \code{ROTO}.
\end{enumerate}
Examples:
\begin{lstlisting}[style=promptstyle]
python scripts/profile_table.py input.csv --run-dir tmp/run
python scripts/audit_v1_8_phase6_app_like_general_regression.py
\end{lstlisting}
For scientific routes:
\begin{lstlisting}[style=promptstyle]
python scripts/audit_v1_8_phase7_app_like_science_astro_regression.py
\end{lstlisting}

\section{Optional, legacy, and domain-specific routes}
\subsection{Optional}
Optional means that an external backend or platform is required, not that
the route is broken. It must not become a core dependency. STILTS, APT,
DuckDB, iWork, Keynote, and TEAREDUCE need preflight and a clear block when
they do not apply.
\subsection{Legacy}
IRAF, fxcor, and iSTARMOD are real but narrow tools. They should not be the
first normal-mode app action. Expert mode may block them or guide the user
toward a safe copy.
\subsection{Domain-specific}
Not every pattern transfers. \code{li6708_equivalent_width_workbench.py measure}
targets a particular spectral line, not general measurement.
\code{catalog_workbench.py crossmatch-sky} transfers tolerance-based matching
but remains sky-only.

\section{Scientific Workbench}
Reference: \path{references/v1-8-scientificworkbench-integration.md}.
Compatibility interpretation:
\begin{itemize}[leftmargin=1.5em]
\item Scientific Workbench invokes local capabilities.
\item AI can plan and explain; local capabilities execute.
\item The app already uses run folders, summaries, manifests, and artifact discovery.
\item v1.8 strengthens the backend to make this integration less fragile.
\end{itemize}
\textbf{v1.8 prepares the skill; v2.0 will synchronize the app.}
This historical guide does not claim that Scientific Workbench consumes
the entire v1.8 contract.

\section{v1.8 regressions}
\begin{longtable}{>{\raggedright\arraybackslash}p{0.55\linewidth}>{\raggedright\arraybackslash}p{0.34\linewidth}}
\toprule
\textbf{Gate} & \textbf{Coverage} \\
\midrule
\path{audit_v1_8_app_ready_contract_regression.py} & JSON contract, examples, and legacy compatibility. \\
\path{audit_v1_8_app_readiness_matrix_regression.py} & 59 rows, counts, and Top 10 P1 priorities. \\
\path{audit_v1_8_phase4_run_bundles_regression.py} & Run bundles and artifact types. \\
\path{audit_v1_8_phase5_router_regression.py} & Dry-run routing on six inputs. \\
\path{audit_v1_8_phase6_app_like_general_regression.py} & Tables, documents, notebooks, ZIP, time series, sensors, and a personal folder. \\
\path{audit_v1_8_phase7_app_like_science_astro_regression.py} & Science, astronomy, optional, and legacy routes. \\
\path{audit_v1_8_phase8_docs_regression.py} & Guide, references, PDF, and key terms. \\
\bottomrule
\end{longtable}

\section{Work reserved for v1.9}
The historical plan for v1.9 is consolidation:
\begin{itemize}[leftmargin=1.5em]
\item Reduce P2 wrapper and artifact-typing debt.
\item Stabilize envelope names and fields.
\item Extend \code{--run-dir} to more priority routes.
\item Strengthen regressions while retaining the same principles.
\item Decide how to handle editable/installed differences before final synchronization.
\item Prepare a stable matrix for app integration.
\end{itemize}

\section{Work reserved for v2.0}
The historical v2.0 scope can be larger:
\begin{itemize}[leftmargin=1.5em]
\item Synchronize Scientific Workbench with the mature skill backend.
\item Consume \code{typed_artifacts}, \code{app_hints}, and \code{next_actions}.
\item Use dry-run routing as a deterministic planning fallback.
\item Separate normal, expert, optional, legacy, and maintainer modes in the UI.
\item Align app run folders with the run bundle contract.
\item Keep provenance, logs, manifests, and the read-only policy visible to users.
\end{itemize}

\section{Historical Phase 8 checklist}
\begin{itemize}[leftmargin=1.5em]
\item v1.8 references present.
\item TeX guide and compiled PDF available.
\item PDF text extracted and validated.
\item No public-surface drift under \code{sync_public_surface_docs.py --check}.
\item Scientific Workbench not edited.
\item Installed skill not synchronized.
\end{itemize}
\end{document}
